\section{Introduction}
\label{sec:intro}

Retrieval-Augmented Generation (RAG)~\cite{lewis2020retrieval} has emerged as the standard architecture for building knowledge-grounded question-answering systems. Rather than relying solely on parametric knowledge stored in model weights, RAG pipelines retrieve relevant documents from an external corpus and condition answer generation on the retrieved context. This approach reduces hallucination~\cite{shuster2021retrieval} and enables domain-specific applications without costly fine-tuning.

A production RAG pipeline involves multiple interacting components: a \emph{chunking strategy} splits source documents into retrievable segments; an \emph{embedding model} maps chunks and queries to a shared vector space; a \emph{retrieval method} selects the most relevant chunks; an optional \emph{reranker} re-orders them; a \emph{prompt template} frames the generation task; and a \emph{generation model} produces the final answer. Each component offers multiple design choices, and the resulting configuration space is combinatorially large. Small changes in chunk size, retrieval strategy, or prompt wording can shift faithfulness scores by more than ten percentage points~\cite{chen2024benchmarking}, yet practitioners often select configurations based on intuition rather than systematic evaluation.

Several evaluation frameworks have been proposed. RAGAS~\cite{es2023ragas} provides automated metrics for answer quality and context usage. ARES~\cite{esphandiary2024areval} offers prediction-powered evaluation. RAGChecker~\cite{izrailovsky2024rag} diagnoses fine-grained retrieval and generation failures. However, these tools evaluate a \emph{fixed} pipeline configuration. They do not systematically sweep the configuration space, compare component interactions, or track experiments for reproducibility. Practitioners are left to orchestrate multi-configuration benchmarks manually, often with ad-hoc scripts that lack structured metric collection.

We address this gap with \toolname{} (Non-Invasive Modular Benchmarking Framework), a framework designed around three principles:

\begin{enumerate}
  \item \textbf{Non-invasive integration.} The framework wraps existing RAG components without restructuring the pipeline. It connects to Ollama~\footnote{\url{https://ollama.com}} and OpenAI-compatible endpoints (vLLM, TGI, LM Studio) through provider routing, requiring only model identifier strings.
  \item \textbf{Systematic configuration exploration.} A declarative \texttt{.env} configuration specifies parameter ranges for all components. The framework generates the cartesian product and benchmarks every combination, producing structured results suitable for direct comparison.
  \item \textbf{Hybrid evaluation.} Each run is assessed through LLM-based quality metrics (faithfulness, context recall, semantic similarity via RAGAS), established IR measures (nDCG@k, Hit@k, Recall@k), NLG metrics (BERTScore, METEOR, ROUGE-L, BLEU), and runtime indicators (time-to-first-token, tokens per second).
\end{enumerate}

We validate \toolname{} by benchmarking 182 configurations across seven generation models spanning local, cloud-hosted, and API deployments on 100-question samples from the SQuAD dataset~\cite{rajpurkar2016squad}. All experiments are tracked via MLflow for full reproducibility.

The main contributions are:
\begin{itemize}
  \item A modular, non-invasive benchmarking framework for systematic RAG configuration analysis.
  \item A hybrid evaluation methodology combining LLM-based, IR, and NLG metrics with runtime measurements.
  \item An empirical study of 182 configurations revealing non-obvious interactions between chunking, retrieval, prompting, and generation components.
\end{itemize}

The remainder of this paper covers related work, framework architecture, experimental setup, results, validity threats, and conclusions.
